\section{Projektvorgehen und Untersuchungsgegenstand}
\label{sec:vorgehen}

\subsection{Arbeitsteilung}

Die Bearbeitung wurde entlang der beiden Themen aufgeteilt:

\begin{table}[H]
  \centering
  \caption{Arbeitsteilung innerhalb der Projektgruppe}
  \begin{tabularx}{\textwidth}{lX}
    \toprule
    \textbf{Thema} & \textbf{Verantwortliche und Aufgaben}\\
    \midrule
    XSS & Arvid Freese und Linus Paliga: Grundlagen, Suche nach
    Quellen und Senken, Analyse der \texttt{vue-cal}-Schwachstelle,
    XSS-Demo und Fix\\
    CORS & Tom Markmann und Raik Ole Kannenberg: Grundlagen, Analyse
    von Spring Security, Vite, Keycloak und Datei-Endpunkten,
    CORS-Demo und sichere Konfiguration\\
    \bottomrule
  \end{tabularx}
\end{table}

Anders als in einem iterativen Issue-Prozess wurden zunächst die
theoretischen Grundlagen getrennt erarbeitet. Anschließend prüften
beide Teilgruppen den Scheduler aus ihrer jeweiligen Perspektive. Die
Ergebnisse wurden in den Markdown-Berichten
\texttt{xss-analyse-und-demo.md} und
\texttt{cors-analyse-und-demo.md} konsolidiert und für die gemeinsame
Präsentation zusammengeführt.

\subsection{Scheduler-Architektur}

Der untersuchte Scheduler verwendet Spring Boot 3.0.1 und Java 17 im
Backend sowie Vue 3 und Vite im Frontend. Die Authentifizierung
erfolgt über Keycloak; die API validiert Bearer-JWTs als OAuth2 Resource Server.

\begin{table}[H]
  \centering
  \caption{Lokale Komponenten und Origins}
  \begin{tabular}{lll}
    \toprule
    \textbf{Komponente} & \textbf{Adresse} & \textbf{Rolle}\\
    \midrule
    Keycloak & \texttt{http://localhost:8080} & Identitätsprovider\\
    Spring-API & \texttt{http://localhost:8079} & Backend und Datenzugriff\\
    Vue-Frontend & \texttt{http://localhost:8078} & reguläre Webanwendung\\
    PostgreSQL & \texttt{localhost:5432} & persistente Daten\\
    Demo-Seite & \texttt{http://localhost:5173} & fremde Origin für CORS-Demo\\
    \bottomrule
  \end{tabular}
\end{table}

Das Frontend ruft die API relativ über \texttt{/api/...} auf.
\texttt{web/vite.config.js} leitet diese Requests lokal an Port 8079
weiter. Aus Sicht des Browsers bleiben die normalen Requests damit
same-origin zu Port 8078.

\subsection{Prüfmethode}

Die Analyse kombinierte statische Quellcodeprüfung,
Konfigurationsvergleich und kontrollierte lokale Demonstrationen. Für
XSS wurden nutzerkontrollierbare Felder bis zu ihren Ausgabestellen
verfolgt und eigene sowie externe HTML-Senken gesucht. Für CORS
wurden Spring Security, Vite-Proxy, Frontend-Requests,
Keycloak-Web-Origins, Nginx und öffentliche Datei-Endpunkte geprüft.

Die weiteren relevanten Entwicklungsbranches wurden ebenfalls auf das
VueCal-Muster geprüft; die ursprüngliche Schwachstelle ist dort
vorhanden. Der dauerhafte Fix wurde ausschließlich auf
\path{fix/vue-cal-xss-safe-rendering} mit Commit \texttt{08567da9}
umgesetzt und geprüft, und zwar ohne Demo-Umschaltung und ohne schwache
CORS-Konfiguration. Er ist bereit für den Merge in
\texttt{capstone-main}. Sämtliche anderen zum Berichtsstand geprüften
Remote-Branches weisen das ursprüngliche anfällige VueCal-Muster
weiterhin auf. Der Demo-Stand
\path{origin/its-xss-cors-security-break} behält bewusst sichere und
unsichere Varianten für den direkten Vergleich.
